\chapter*{Mở đầu}
\addcontentsline{toc}{chapter}{Mở đầu}
\setcounter{page}{1}
\pagenumbering{arabic}
Ngày nay, với tốc độ tăng trưởng mạnh mẽ của công nghệ thông tin, xu hướng phát triển phần mềm đang đổi mới với tốc độ không ngừng. Những khảo sát gần đây cho thấy phạm vi và độ phức tạp của các hệ thống phần mềm đã tăng lên đáng kể. Các công ty luôn mong muốn cải tiến quy trình phát triển để tối ưu chi phí, cân đối nguồn lực cần đầu tư để tạo mới, mở rộng hay bảo trì các sản phẩm phần mềm. Tái sử dụng phần mềm là giải pháp hợp lý nhất cho vấn đề này. Cụ thể, tái sử dụng là khái niệm phát triển phần mềm bằng các tài nguyên có hiện có thay vì xây dựng mọi thứ từ đầu \cite{Krueger}. Tất cả các thành phần trong quá trình phát triển phần mềm đều có thể được chuyển đổi thành các tài nguyên phục vụ sử dụng lại, bao gồm mã nguồn, kiến trúc, đặc tả, thiết kế, hay kể cả tài liệu. \par
Một xu hướng mới cũng đang được các công ty phần mềm quan tâm là việc tập trung nguồn lực để phát triển theo dòng sản phẩm (SPL\footnote{Software Product Lines}) thay vì từng sản phẩm riêng biệt \cite{Guendouz2014}. Những hệ thống dạng này thường mang ưu thế về sự linh hoạt giúp đáp ứng nhiều yêu cầu khác nhau của khách hàng. Cùng với nguyên tắc tái sử dụng tối đa, công đoạn thiết kế và sản xuất một sản phẩm mới cũng được rút ngắn đáng kể. So với các phương thức truyền thống, quá trình thiết kế hệ thống dòng sản phẩm tương đối khác biệt khi tập trung phân tích trước bộ khung tổng thể và cài đặt mã nguồn cho từng tính năng độc lập. Với từng yêu cầu đầu vào, đặc tả về bộ chức năng của sản phẩm hoàn chỉnh sẽ được mô hình hóa dưới dạng cấu hình. Một cấu hình bao gồm tất cả các tính năng sẵn có trong hệ thống kèm theo trạng thái bật hay tắt, biểu thị rằng tính năng đó có được đưa vào sản phẩm cuối cùng hay không. Với cơ chế như vậy, mã nguồn của chương trình là kết quả của quá trình biên soạn, kết hợp các tính năng theo từng cấu hình tương ứng. \par
Về vấn đề đảm bảo chất lượng cho hệ thống SPL, những nghiên cứu gần đây đang bắt đầu chú trọng hơn về việc tối ưu quy trình kiểm thử cho dòng sản phẩm. Mặc dù đã có nhiều phương pháp được đề xuất giúp giải quyết từng bài toán trong lĩnh vực này \cite{Piwowarski257635, Namin1572280,Kim961347, Lamancha16573, ARRIETA201818}, cách tiến hành đánh giá độ hiệu quả của chúng còn gặp nhiều hạn chế bởi sự thiếu hụt dữ liệu trong quá trình thử nghiệm. Trong thực tế, các dự án về dòng sản phẩm thường ít khi được công khai, nếu có thì thường chỉ bao gồm mã nguồn đơn thuần thay vì chứa cả dữ liệu thử được sinh trực tiếp trên những hệ thống ấy. Đã có những nghiên cứu đi sâu để khắc phục vấn đề về dữ liệu bằng cách xây dựng bộ công cụ hỗ trợ tạo dữ liệu đánh giá\footnote{Benchmark Dataset}, với mục tiêu để so sánh khả năng hoạt động các kĩ thuật kiểm thử trong hệ thống SPL với nhau \cite{Herrejon2014,Fischer3167350}. Tuy nhiên những bộ dữ liệu được sinh ra theo phương thức trước đây chỉ dừng lại ở mức xét tập cấu hình thay vì đi sâu vào mã nguồn để tìm hiểu nguyên nhân gây lỗi. Ngoài ra trong quá trình thực thi, các thông tin về độ phủ dòng lệnh cũng ít khi được thu thập lại, mặc dù đây được coi như những dữ kiện bắt buộc để làm đầu vào cho các kĩ thuật khoanh vùng, sửa lỗi của chương trình.\par
Từ những khó khăn như vậy, khóa luận này sẽ trình bày một giải pháp giúp tạo ra tập dữ liệu đánh giá bao gồm các hệ thống dòng sản phẩm có lỗi kèm theo bộ kịch bản kiểm thử đầy đủ. Trong quá trình thực thi, những thông tin về vị trí gây lỗi và độ phủ từng dòng lệnh cũng sẽ được ghi lại. Từ bộ dữ liệu được sinh, các kĩ thuật kiểm thử và khoanh vùng lỗi được đề xuất trước đây sẽ được đưa ra để so sánh về độ hữu hiệu theo cách hoàn toàn khách quan.

Khóa luận sẽ được trình bày theo cấu trúc như sau. Chương 1 sẽ giới thiệu tổng quan về hệ thống dòng sản phẩm với quy trình xây dựng và cấu trúc các bộ phận liên quan. Chương 2 sẽ nhắc lại sơ lược một số kiến thức về kiểm thử phần mềm và đưa ra những phương pháp để kiểm thử cho một dòng sản phẩm cụ thể. Giải pháp được đề xuất trong khóa luận sẽ được mô tả chi tiết ở Chương 3. Tiếp theo, Chương 4 tập trung vào việc thử nghiệm phương pháp và báo cáo kết quả chi tiết. Phần cuối cùng rút ra kết luận và hướng mở rộng trong tương lai.

\cleardoublepage